\[
  F = -\frac{GM_{BH}M_*}{R^2} = -\frac{M_*v^2}{R}
\]
\hfill(2 Point)

But $v = \dfrac{2\pi R}{P}$. Therefore
\[
  \frac{GM_{BH}}{R^2} = 4\pi^2\,\frac{R}{P^2}
  \qquad\text{or}\qquad
  GM_{BH} = 4\pi^2\,\frac{R^3}{P^2}
\]
\hfill(2 Points)

Similarly:
\[
  GM_\odot = 4\pi^2\,\frac{(1\,AU)^3}{(1\,yr)^2}
  \qquad\text{or}\qquad
  G = \frac{1}{M_\odot}\,4\pi^2\,\frac{(1\,AU)^3}{(1\,yr)^2}
\]
\hfill(2 Point)

From Kepler's 3rd law we get:
\[
  \frac{M_{BH}}{M_\odot} = \frac{(R/1\,AU)^3}{(P/1\,yr)^2}
\]
\hfill(1 Point)

Inserting the given data we find the distance of the star from the black hole:
\[
  R = \frac{0.12}{200{,}000}\,(8000)(3\times10^{18}\,\mathrm{cm})
    = 1.4\times10^{16}\,\mathrm{cm} = 960\,AU
\]
\hfill(2 Point)

Therefore:
\[
  \frac{M_{BH}}{M_\odot} = \frac{(960)^3}{(15)^2} = 4\times10^{6}
\]
form which % sic
we calculate the mass of the black hole:
\[
  M_{BH} = 4\times10^{6}\,M_\odot
\]
\hfill(1 Point)
